% ─────────────────────────────────────────────────────────────────────────────
%  PART 2 SECTION DIVIDER
% ─────────────────────────────────────────────────────────────────────────────
\sectionframe{2}{Mechanistic Interpretability}%
  {Paper \textcircled{2} — Locate, Steer \& Improve}%
  {Zhang et al.\,(2026) $\cdot$ arXiv 2601.14004}

% ═════════════════════════════════════════════════════════════════════════════
%  SLIDE 18 — WHAT IS MI?
% ═════════════════════════════════════════════════════════════════════════════
\begin{frame}{What Is Mechanistic Interpretability?}
  \begin{ucrinfo}
    \small\itshape
    ``Mechanistic Interpretability (MI) aims to reverse-engineer neural networks —
    decomposing their intricate computations into understandable components
    and causal mechanisms.'' \hfill\muted{\footnotesize Zhang et al., 2026}
  \end{ucrinfo}
  \vspace{0.4em}

  \begin{columns}[T]
    \column{0.48\textwidth}
    \begin{tcolorbox}[enhanced, colback=UCRlgray, colframe=UCRmgray,
        sharp corners, fonttitle=\bfseries\small\color{UCRmgray},
        title=Traditional Interpretability]
      \begin{itemize}\footnotesize\setlength\itemsep{3pt}
        \item[\color{UCRmgray}$\circ$] SHAP, LIME, attention visualization
        \item[\color{UCRmgray}$\circ$] Explains \emph{what}, not \emph{how}
        \item[\color{UCRmgray}$\circ$] Treats model as a black box
        \item[\color{UCRmgray}$\circ$] Limited ability to enable targeted interventions
      \end{itemize}
    \end{tcolorbox}

    \column{0.48\textwidth}
    \begin{tcolorbox}[enhanced, colback=UCRlgray, colframe=UCRteal,
        toprule=2pt, sharp corners, fonttitle=\bfseries\small\color{UCRteal},
        title=Mechanistic Interpretability]
      \begin{itemize}\footnotesize\setlength\itemsep{3pt}
        \item Identifies neurons, heads, circuits responsible for behaviors
        \item Reveals causal mechanisms inside the model
        \item Enables surgical interventions without full retraining
        \item Translates insights into practical model improvements
      \end{itemize}
    \end{tcolorbox}
  \end{columns}
  \vspace{0.3em}
  \begin{ucrinfo}[UCRteal]
    \footnotesize\bfseries Core shift:\normalfont\ from observational science
    (`what happens?') $\to$ actionable intervention discipline (`how do we fix it?')
  \end{ucrinfo}
\end{frame}

% ═════════════════════════════════════════════════════════════════════════════
%  SLIDE 19 — LOCATE $\to$ STEER $\to$ IMPROVE PIPELINE
% ═════════════════════════════════════════════════════════════════════════════
\begin{frame}{The ``Locate $\to$ Steer $\to$ Improve'' Pipeline \papertag[UCRteal]{Paper \textcircled{2}}}
  \vspace*{0.3cm}
  \begin{tikzpicture}[
    block/.style={rectangle, rounded corners=2pt, draw=#1, fill=UCRlgray,
                  text width=3.5cm, minimum height=4.5cm, text centered,
                  align=center},
    arr/.style={-{Stealth[length=6pt]}, UCRmgray, very thick},
  ]
    \node[block=UCRteal, label={[font=\Huge\bfseries\color{UCRteal},
          opacity=0.2]center:1}] (l) at (0,0)
      {{\bfseries\color{UCRteal}\large LOCATE}\\[0.15em]
       \footnotesize\itshape(Diagnose)\\[0.5em]
       \footnotesize Identify which components — neurons, heads, circuits,
       SAE features — are causally responsible for a specific behavior.};

    \node[block=UCRnavy, label={[font=\Huge\bfseries\color{UCRnavy},
          opacity=0.2]center:2}] (s) at (4.1,0)
      {{\bfseries\color{UCRnavy}\large STEER}\\[0.15em]
       \footnotesize\itshape(Intervene)\\[0.5em]
       \footnotesize Actively manipulate identified components: amplitude
       manipulation, targeted optimization, vector arithmetic.};

    \node[block=UCRgreen, label={[font=\Huge\bfseries\color{UCRgreen},
          opacity=0.2]center:3}] (i) at (8.2,0)
      {{\bfseries\color{UCRgreen}\large IMPROVE}\\[0.15em]
       \footnotesize\itshape(Apply)\\[0.5em]
       \footnotesize Translate insights into practical improvements in
       alignment, capability, and efficiency — without full retraining.};

    \draw[arr] (l.east) -- (s.west);
    \draw[arr] (s.east) -- (i.west);
  \end{tikzpicture}
  \vspace{0.3em}
  \centerline{\muted{\footnotesize\itshape
    Applicable Objects: Token Embedding $\cdot$ Residual Stream $\cdot$ MHA $\cdot$ FFN $\cdot$ SAE Features}}
\end{frame}

% ═════════════════════════════════════════════════════════════════════════════
%  SLIDE 20 — CORE INTERPRETABLE OBJECTS OVERVIEW
% ═════════════════════════════════════════════════════════════════════════════
\begin{frame}{Core Interpretable Objects of LLMs \papertag[UCRteal]{Paper \textcircled{2}}}
  {\small The `interpretable objects' are the units MI methods analyze and
    intervene upon inside a Transformer LLM:}
  \vspace{0.3em}

  \begin{tcolorbox}[enhanced, colback=UCRlgray, colframe=UCRmgray,
      borderline west={3pt}{0pt}{UCRteal}, sharp corners, top=3pt, bottom=3pt]
    \small\bfseries\tealc{1. Token Embedding}\normalfont\
    — Maps discrete tokens to continuous vectors; entry point; encodes lexical \& positional meaning.
  \end{tcolorbox}\vspace{2pt}
  \begin{tcolorbox}[enhanced, colback=UCRlgray, colframe=UCRmgray,
      borderline west={3pt}{0pt}{UCRnavy}, sharp corners, top=3pt, bottom=3pt]
    \small\bfseries\navy{2. Transformer Block \& Residual Stream}\normalfont\
    — Central `highway' of information. Features are linear combinations of all prior component outputs.
  \end{tcolorbox}\vspace{2pt}
  \begin{tcolorbox}[enhanced, colback=UCRlgray, colframe=UCRmgray,
      borderline west={3pt}{0pt}{UCRpurple}, sharp corners, top=3pt, bottom=3pt]
    \small\bfseries{\color{UCRpurple}3. Multi-Head Attention (MHA)}\normalfont\
    — Manages information routing; key units: QK (where to attend), OV (what to copy).
  \end{tcolorbox}\vspace{2pt}
  \begin{tcolorbox}[enhanced, colback=UCRlgray, colframe=UCRmgray,
      borderline west={3pt}{0pt}{UCRorange}, sharp corners, top=3pt, bottom=3pt]
    \small\bfseries{\color{UCRorange}4. Feed-Forward Network (FFN)}\normalfont\
    — Acts as key-value memory; stores factual associations; neurons encode specific world-knowledge tokens.
  \end{tcolorbox}\vspace{2pt}
  \begin{tcolorbox}[enhanced, colback=UCRlgray, colframe=UCRmgray,
      borderline west={3pt}{0pt}{UCRgreen}, sharp corners, top=3pt, bottom=3pt]
    \small\bfseries\greenc{5. Sparse Autoencoder (SAE) Feature}\normalfont\
    — Disentangles polysemantic neurons into monosemantic features via sparse coding — a `microscope' for LLM representations.
  \end{tcolorbox}
\end{frame}

% ═════════════════════════════════════════════════════════════════════════════
%  SLIDE 21 — TOKEN EMBEDDING & RESIDUAL STREAM
% ═════════════════════════════════════════════════════════════════════════════
\begin{frame}{Token Embedding \& Residual Stream}
  \begin{columns}[T]
    \column{0.48\textwidth}
    {\small\bfseries\tealc{Token Embedding}}
    \begin{itemize}\footnotesize\setlength\itemsep{3pt}
      \item $W_E \in \mathbb{R}^{|V|\times d_{\text{model}}}$ maps token ids to vectors
      \item Initial residual stream state:
        $x^0_i = W_E[t_i] + p_i$
      \item Modern LLMs use RoPE (Rotary Position Embedding) —
        positional info applied to Q/K, not the stream
      \item \textbf{Logit Lens}: project intermediate states back to
        vocabulary to inspect per-layer predictions
    \end{itemize}

    \column{0.48\textwidth}
    {\small\bfseries\navy{Residual Stream}}
    \begin{itemize}\footnotesize\setlength\itemsep{3pt}
      \item Central communication channel across all layers
      \item Update dynamics:
        $x^{l+1} = x^l + \text{MHA}(x^l) + \text{FFN}(x^l + \text{MHA}(x^l))$
      \item Additive structure: final prediction is a linear combination
        of all component contributions
      \item Enables: Logit Lens, causal mediation analysis, path patching
    \end{itemize}
  \end{columns}
  \vspace{0.4em}
  \begin{tikzpicture}
    \foreach \i/\lbl/\col in {
      0/Input Tokens/UCRteal, 1/Embed ($x^0$)/UCRnavy,
      2/Block 1/UCRnavy, 3/Block $L$/UCRnavy, 4/Output Logits/UCRgreen}
    {
      \node[rectangle, fill=\col, text=white, font=\tiny\bfseries,
            minimum width=1.8cm, minimum height=0.65cm, text centered]
        at (\i*2.2, 0) {\lbl};
      \ifnum\i<4
        \draw[-{Stealth[length=4pt]}, UCRmgray] (\i*2.2+0.9, 0) -- (\i*2.2+1.3, 0);
      \fi
    }
  \end{tikzpicture}
  \vspace{0.1em}
  \centerline{\muted{\tiny\itshape
    Residual stream flows through all layers — each block reads from it and writes back additively}}
\end{frame}

% ═════════════════════════════════════════════════════════════════════════════
%  SLIDE 22 — MHA & FFN
% ═════════════════════════════════════════════════════════════════════════════
\begin{frame}{Multi-Head Attention \& Feed-Forward Networks}
  \begin{columns}[T]
    \column{0.48\textwidth}
    {\small\bfseries\tealc{Multi-Head Attention (MHA)}}
    \begin{itemize}\footnotesize\setlength\itemsep{3pt}
      \item $H$ independent heads per layer; each routes information between positions
      \item \textbf{QK unit}: determines \emph{where} to attend (attention pattern $A^{l,h}$)
      \item \textbf{OV unit}: determines \emph{what} information to copy
      \item Heads specialize: `induction heads', `name mover heads', `copy heads'
      \item Causal attribution identifies heads mediating specific behaviors
        (refusal, factual recall)
    \end{itemize}

    \column{0.48\textwidth}
    {\small\bfseries\navy{Feed-Forward Networks (FFN)}}
    \begin{itemize}\footnotesize\setlength\itemsep{3pt}
      \item FFN = key-value memory (Geva et al., 2021)
      \item Keys ($W_{\text{in}}$ rows): recognize input patterns
      \item Values ($W_{\text{out}}$ rows): project knowledge onto residual stream
      \item Individual neurons encode specific world knowledge
        (e.g., `Paris is capital of France')
      \item Factual editing methods (ROME, MEMIT) directly update FFN weights
    \end{itemize}
  \end{columns}
  \vspace{0.3em}
  \begin{ucrwarn}
    \small\bfseries Superposition Problem:\normalfont\ Neurons are
    \emph{polysemantic} — a single neuron activates for multiple unrelated
    concepts. This is the core challenge MI must solve, and why SAEs are needed.
  \end{ucrwarn}
\end{frame}

% ═════════════════════════════════════════════════════════════════════════════
%  SLIDE 23 — SPARSE AUTOENCODERS
% ═════════════════════════════════════════════════════════════════════════════
\begin{frame}{Sparse Autoencoders — Resolving Superposition}
  \begin{columns}[T]
    \column{0.52\textwidth}
    {\small\bfseries\redc{The Problem: Superposition}}\\[0.2em]
    {\footnotesize Neural networks store more features than neurons by encoding
      them as nearly-orthogonal directions in high-dimensional space.
      One neuron $\to$ many unrelated concepts.}
    \vspace{0.4em}

    {\small\bfseries\tealc{The SAE Solution}}
    \begin{itemize}\footnotesize\setlength\itemsep{3pt}
      \item SAE = encoder-decoder trained on frozen LLM activations
      \item Projects dense activations $\to$ high-dim sparse latent space
      \item Expansion factor: $d_{\text{SAE}} = 16\times\text{–}128\times\, d_{\text{model}}$
      \item Training loss: $\mathcal{L} = \|x^l - \hat{x}^l\|_2^2 + \lambda\|a\|_1$
      \item Each SAE feature $f_j$ = one interpretable semantic direction
      \item Templeton et al.\ (2024): features include `harmful content',
        `toxicity', `refusal'
    \end{itemize}

    \column{0.44\textwidth}
    \begin{tcolorbox}[enhanced, colback=UCRlgray, colframe=UCRmgray,
        sharp corners, fonttitle=\bfseries\small\color{UCRteal},
        title=SAE Architecture]
      \begin{tikzpicture}
        \node[rectangle, fill=UCRnavy, text=white, font=\tiny\bfseries,
              minimum width=3.5cm, minimum height=0.6cm, text centered]
          (dense) {Dense Activation $x^l$\;($d_{\text{model}}$)};
        \draw[-{Stealth[length=4pt]}, UCRmgray]
          (dense.south) -- +(0,-0.3);
        \node[rectangle, fill=UCRteal, text=white, font=\tiny\bfseries,
              minimum width=3.5cm, minimum height=0.6cm, text centered,
              below=0.3cm of dense]
          (enc) {SAE Encoder $\to$ Sparse $a$\;($d_{\text{SAE}}$)};
        \draw[-{Stealth[length=4pt]}, UCRmgray]
          (enc.south) -- +(0,-0.3);
        \node[rectangle, fill=UCRnavy, text=white, font=\tiny\bfseries,
              minimum width=3.5cm, minimum height=0.6cm, text centered,
              below=0.3cm of enc]
          (dec) {SAE Decoder $\to$ Reconstruction $\hat{x}^l$};
      \end{tikzpicture}
      \vspace{0.3em}
      \scriptsize
      $d_{\text{SAE}} \gg d_{\text{model}}$\\
      Each feature $a_j$ is monosemantic\\
      Dead latent problem: some features never activate
    \end{tcolorbox}
  \end{columns}
\end{frame}

% ═════════════════════════════════════════════════════════════════════════════
%  SLIDE 24 — LOCALIZING METHODS
% ═════════════════════════════════════════════════════════════════════════════
\begin{frame}{Localizing Methods — Finding the Right Components}
  {\footnotesize\itshape Goal: identify internal objects causally responsible
    for a specific behavior.}
  \vspace{0.2em}

  \begin{columns}[T]
    \column{0.48\textwidth}
    \begin{tcolorbox}[enhanced, colback=UCRlgray, colframe=UCRmgray,
        borderline west={3pt}{0pt}{UCRteal}, sharp corners,
        fonttitle=\bfseries\footnotesize\color{UCRteal},
        title=Magnitude Analysis, top=3pt, bottom=3pt]
      \scriptsize Rank objects by activation magnitude or weight norms.
      High-magnitude elements $\approx$ salient components.
    \end{tcolorbox}\vspace{3pt}
    \begin{tcolorbox}[enhanced, colback=UCRlgray, colframe=UCRmgray,
        borderline west={3pt}{0pt}{UCRnavy}, sharp corners,
        fonttitle=\bfseries\footnotesize\color{UCRnavy},
        title=Causal Attribution, top=3pt, bottom=3pt]
      \scriptsize Activation patching: replace one activation from reference run;
      measure causal effect on output. Precise causal identification.
    \end{tcolorbox}\vspace{3pt}
    \begin{tcolorbox}[enhanced, colback=UCRlgray, colframe=UCRmgray,
        borderline west={3pt}{0pt}{UCRpurple}, sharp corners,
        fonttitle=\bfseries\footnotesize\color{UCRpurple},
        title=Gradient Detection, top=3pt, bottom=3pt]
      \scriptsize Integrated Gradients, SmoothGrad, saliency maps.
      Identifies which inputs drive which computations.
    \end{tcolorbox}

    \column{0.48\textwidth}
    \begin{tcolorbox}[enhanced, colback=UCRlgray, colframe=UCRmgray,
        borderline west={3pt}{0pt}{UCRorange}, sharp corners,
        fonttitle=\bfseries\footnotesize\color{UCRorange},
        title=Probing, top=3pt, bottom=3pt]
      \scriptsize Train lightweight classifiers on internal activations.
      Tests what information is encoded where.
    \end{tcolorbox}\vspace{3pt}
    \begin{tcolorbox}[enhanced, colback=UCRlgray, colframe=UCRmgray,
        borderline west={3pt}{0pt}{UCRgreen}, sharp corners,
        fonttitle=\bfseries\footnotesize\color{UCRgreen},
        title=Vocabulary Projection (Logit Lens), top=3pt, bottom=3pt]
      \scriptsize Project residual stream states to vocabulary space.
      Reveals what each layer `predicts' at each step.
    \end{tcolorbox}\vspace{3pt}
    \begin{tcolorbox}[enhanced, colback=UCRlgray, colframe=UCRmgray,
        borderline west={3pt}{0pt}{UCRred}, sharp corners,
        fonttitle=\bfseries\footnotesize\color{UCRred},
        title=Circuit Discovery, top=3pt, bottom=3pt]
      \scriptsize Extract minimal subgraphs (circuits) implementing a task.
      Path patching traces information flow.
    \end{tcolorbox}
  \end{columns}
\end{frame}

% ═════════════════════════════════════════════════════════════════════════════
%  SLIDE 25 — CAUSAL ATTRIBUTION & CIRCUIT DISCOVERY
% ═════════════════════════════════════════════════════════════════════════════
\begin{frame}{Deep Dive: Causal Attribution \& Circuit Discovery}
  \begin{columns}[T]
    \column{0.48\textwidth}
    {\small\bfseries\navy{Causal Attribution via Activation Patching}}
    \begin{enumerate}\footnotesize\setlength\itemsep{3pt}
      \item Run model on `clean' and `corrupted' prompt
      \item Patch one activation at a time: clean $\to$ corrupted run
      \item Measure logit difference $\Delta L$
      \item Components with large $\Delta L$ are causally important
    \end{enumerate}
    \vspace{0.2em}
    {\footnotesize ROME used causal tracing to locate factual memories
      in specific MLP layers. Path patching (extension) traces information
      flow along specific edges.}

    \column{0.48\textwidth}
    {\small\bfseries\tealc{Circuit Discovery}}
    \begin{itemize}\footnotesize\setlength\itemsep{3pt}
      \item Circuit = minimal component subgraph implementing a task
      \item Nodes = heads/MLPs; edges = information flow
      \item Method: iterative edge pruning
      \item Famous examples: IOI circuit; Greater-Than circuit
      \item Circuit Tracing (Anthropic 2025): computational graphs in Claude
      \item Fix the circuit $\to$ fix the behavior
    \end{itemize}
  \end{columns}
  \vspace{0.3em}
  \begin{ucrinfo}[UCRteal]
    \footnotesize\bfseries Mechanistic Insight:\normalfont\ A small fraction of
    components (sometimes \textless5\% of parameters) is responsible for specific
    capabilities. MI allows `surgical' edits without degrading the rest.
  \end{ucrinfo}
\end{frame}

% ═════════════════════════════════════════════════════════════════════════════
%  SLIDE 26 — STEERING METHODS
% ═════════════════════════════════════════════════════════════════════════════
\begin{frame}{Steering Methods — Intervening on LLM Internals}
  \begin{tcolorbox}[enhanced, colback=UCRlgray, colframe=UCRmgray,
      borderline west={4pt}{0pt}{UCRteal}, sharp corners,
      fonttitle=\bfseries\small\color{UCRteal}, title=Amplitude Manipulation,
      top=3pt, bottom=3pt]
    \footnotesize Inference-time activation control: \textbf{ablation} (zero out),
    \textbf{patching} (replace), \textbf{scaling} (amplify/attenuate).
    Reversible; no weight update. Tang et al.\ (2024): zero out Chinese neurons $\to$
    force English output.
    \hfill\redc{\tiny \textbf{!} fragile if target is polysemantic}
  \end{tcolorbox}
  \vspace{3pt}
  \begin{tcolorbox}[enhanced, colback=UCRlgray, colframe=UCRmgray,
      borderline west={4pt}{0pt}{UCRnavy}, sharp corners,
      fonttitle=\bfseries\small\color{UCRnavy}, title=Targeted Optimization,
      top=3pt, bottom=3pt]
    \footnotesize Localized parameter updates restricted to mask $M$ of selected components:
    $\theta' \leftarrow \theta + (M \odot \Delta\theta)$.
    ROME/MEMIT edit factual knowledge in FFN layers.
    Persistent across inference calls.
    \hfill\redc{\tiny \textbf{!} careful mask selection required}
  \end{tcolorbox}
  \vspace{3pt}
  \begin{tcolorbox}[enhanced, colback=UCRlgray, colframe=UCRmgray,
      borderline west={4pt}{0pt}{UCRgreen}, sharp corners,
      fonttitle=\bfseries\small\color{UCRgreen}, title=Vector Arithmetic,
      top=3pt, bottom=3pt]
    \footnotesize Operate on the residual stream: add/subtract steering vectors
    for semantic directions. Refusal direction: add refusal vector $\to$ suppress
    jailbreaks (Arditi et al., 2024). Representation Engineering (RepE) steers
    via contrastive pairs.
    \hfill\redc{\tiny \textbf{!} linear assumption may break down}
  \end{tcolorbox}
\end{frame}

% ═════════════════════════════════════════════════════════════════════════════
%  SLIDE 27 — APPLICATIONS: ALIGNMENT
% ═════════════════════════════════════════════════════════════════════════════
\begin{frame}{MI Applications — Improve Alignment \papertag[UCRteal]{Paper \textcircled{2}}}
  \begin{tcolorbox}[enhanced, colback=UCRlgray, colframe=UCRmgray,
      borderline west={4pt}{0pt}{UCRgreen}, sharp corners,
      fonttitle=\bfseries\small\color{UCRgreen}, title=Safety \& Reliability,
      top=3pt, bottom=3pt]
    \begin{itemize}\footnotesize\setlength\itemsep{2pt}
      \item Locate safety-critical attention heads/neurons via causal attribution
      \item SAE features for toxicity: suppress at inference time (Goyal et al., 2025)
      \item Refusal direction in residual stream: jailbreaks suppress this direction
        (Arditi et al., 2024)
      \item \textbf{Refusal cliff} (Yin et al., 2025): refusal intent maintained in
        reasoning but suppressed at final generation by a small set of attention heads
    \end{itemize}
  \end{tcolorbox}
  \vspace{3pt}
  \begin{tcolorbox}[enhanced, colback=UCRlgray, colframe=UCRmgray,
      borderline west={4pt}{0pt}{UCRteal}, sharp corners,
      fonttitle=\bfseries\small\color{UCRteal}, title=Fairness \& Bias,
      top=3pt, bottom=3pt]
    \begin{itemize}\footnotesize\setlength\itemsep{2pt}
      \item Identify gender/demographic bias in specific attention heads
      \item Ahsan et al.\ (2025): `Male patch' intervention changes clinical predictions
        for depression — reveals causal link between demographic representations and
        model decisions
    \end{itemize}
  \end{tcolorbox}
  \vspace{3pt}
  \begin{tcolorbox}[enhanced, colback=UCRlgray, colframe=UCRmgray,
      borderline west={4pt}{0pt}{UCRnavy}, sharp corners,
      fonttitle=\bfseries\small\color{UCRnavy}, title=Persona \& Role,
      top=3pt, bottom=3pt]
    \footnotesize Steer model persona via vector arithmetic; representation
    engineering applies contrastive pairs for desired personality traits.
  \end{tcolorbox}
\end{frame}

% ═════════════════════════════════════════════════════════════════════════════
%  SLIDE 28 — APPLICATIONS: CAPABILITY & EFFICIENCY
% ═════════════════════════════════════════════════════════════════════════════
\begin{frame}{MI Applications — Capability \& Efficiency \papertag[UCRteal]{Paper \textcircled{2}}}
  \begin{columns}[T]
    \column{0.48\textwidth}
    {\small\bfseries\tealc{Improve Capability}}

    \begin{tcolorbox}[enhanced, colback=UCRlgray, colframe=UCRmgray,
        borderline west={3pt}{0pt}{UCRteal}, sharp corners,
        fonttitle=\bfseries\footnotesize\color{UCRteal},
        title=Multilingualism, top=3pt, bottom=3pt]
      \scriptsize Locate language-specific neurons; cross-lingual transfer by
      enabling/disabling circuits. Tang et al.\ (2024): zero out Chinese neurons $\to$ force English.
    \end{tcolorbox}\vspace{2pt}
    \begin{tcolorbox}[enhanced, colback=UCRlgray, colframe=UCRmgray,
        borderline west={3pt}{0pt}{UCRnavy}, sharp corners,
        fonttitle=\bfseries\footnotesize\color{UCRnavy},
        title=Knowledge Editing, top=3pt, bottom=3pt]
      \scriptsize Factual knowledge in FFN layers. ROME/MEMIT: targeted
      optimization updates/erases specific facts without full retraining.
    \end{tcolorbox}\vspace{2pt}
    \begin{tcolorbox}[enhanced, colback=UCRlgray, colframe=UCRmgray,
        borderline west={3pt}{0pt}{UCRpurple}, sharp corners,
        fonttitle=\bfseries\footnotesize\color{UCRpurple},
        title=Logic \& Reasoning, top=3pt, bottom=3pt]
      \scriptsize SAE features for reasoning (Galichin et al., 2025); scaling
      reflection features increases reasoning depth; arithmetic circuits discovered.
    \end{tcolorbox}

    \column{0.48\textwidth}
    {\small\bfseries\navy{Improve Efficiency}}

    \begin{tcolorbox}[enhanced, colback=UCRlgray, colframe=UCRmgray,
        borderline west={3pt}{0pt}{UCRnavy}, sharp corners,
        fonttitle=\bfseries\footnotesize\color{UCRnavy},
        title=Pruning \& Layer Skipping, top=3pt, bottom=3pt]
      \scriptsize Magnitude analysis finds redundant heads/layers.
      Men et al.\ (2025): remove low-causal-effect layers $\to$ faster inference.
    \end{tcolorbox}\vspace{2pt}
    \begin{tcolorbox}[enhanced, colback=UCRlgray, colframe=UCRmgray,
        borderline west={3pt}{0pt}{UCRteal}, sharp corners,
        fonttitle=\bfseries\footnotesize\color{UCRteal},
        title=Attention Head Pruning, top=3pt, bottom=3pt]
      \scriptsize Many heads are functionally redundant; retain only
      task-critical heads $\to$ faster inference with minimal quality loss.
    \end{tcolorbox}\vspace{2pt}
    \begin{tcolorbox}[enhanced, colback=UCRlgray, colframe=UCRmgray,
        borderline west={3pt}{0pt}{UCRgreen}, sharp corners,
        fonttitle=\bfseries\footnotesize\color{UCRgreen},
        title=MI-Informed Fine-tuning, top=3pt, bottom=3pt]
      \scriptsize Target LoRA updates to MI-localized components $\to$
      parameter-efficient adaptation with stronger guarantees.
    \end{tcolorbox}
  \end{columns}
\end{frame}

